\documentclass[11pt,a4paper]{article}
\usepackage[T1]{fontenc}
\usepackage[utf8]{inputenc}
\usepackage[a4paper,margin=25mm,headheight=15pt]{geometry}
\usepackage{amsmath,amssymb,amsthm,mathtools}
\usepackage{newpxtext,newpxmath}
\usepackage{microtype,booktabs,tabularx,array,longtable}
\usepackage[dvipsnames]{xcolor}
\usepackage{enumitem,fancyhdr,titlesec,needspace}
\usepackage{url}
\usepackage[colorlinks=true,linkcolor=MidnightBlue,citecolor=MidnightBlue,urlcolor=MidnightBlue,breaklinks=true]{hyperref}
\definecolor{Ink}{HTML}{243747}
\definecolor{Accent}{HTML}{566348}
\titleformat{\section}{\Large\bfseries\color{Ink}}{\thesection}{.65em}{}
\titleformat{\subsection}{\large\bfseries\color{Accent}}{\thesubsection}{.65em}{}
\pagestyle{fancy}\fancyhf{}
\fancyhead[L]{\small Tail-spans and differential transcendence}
\fancyhead[R]{\small Surreal and surcomplex fields}
\fancyfoot[C]{\thepage}
\setlength{\parskip}{.23em}
\setlength{\parindent}{1.05em}
\setlist{itemsep=.12em,topsep=.3em,leftmargin=2em}
\setcounter{tocdepth}{2}
\emergencystretch=2em
\numberwithin{equation}{section}
\newtheorem{theorem}{Theorem}[section]
\newtheorem{lemma}[theorem]{Lemma}
\newtheorem{proposition}[theorem]{Proposition}
\newtheorem{corollary}[theorem]{Corollary}
\theoremstyle{definition}
\newtheorem{definition}[theorem]{Definition}
\newtheorem{example}[theorem]{Example}
\theoremstyle{remark}
\newtheorem{remark}[theorem]{Remark}
\newtheorem{warning}[theorem]{Warning}
\newcommand{\No}{\mathbf{No}}
\newcommand{\SC}{\mathbf{No}[i]}
\newcommand{\R}{\mathbb R}
\newcommand{\C}{\mathbb C}
\newcommand{\Q}{\mathbb Q}
\newcommand{\N}{\mathbb N}
\newcommand{\Z}{\mathbb Z}
\newcommand{\AR}{\mathbb A_{\R}}
\newcommand{\OO}{\mathcal O}
\newcommand{\supp}{\operatorname{supp}}
\newcommand{\Span}{\operatorname{span}}
\newcommand{\trdeg}{\operatorname{trdeg}}
\newcommand{\dtrdeg}{\operatorname{dtrdeg}}
\newcommand{\Gal}{\operatorname{Gal}}
\newcommand{\Frac}{\operatorname{Frac}}
\newcommand{\Spec}{\operatorname{Spec}}
\newcommand{\rank}{\operatorname{rank}}
\newcommand{\BM}{D_{\mathrm{BM}}}
\newcommand{\EDer}{\delta}
\newcommand{\base}{B}
\newcommand{\cf}{\mathfrak c}
\newcommand{\ff}[2]{(#1)_{\underline{#2}}}
\newcommand{\rf}[2]{(#1)^{\overline{#2}}}
\newcommand{\repocommit}{608dd23c0539fce73723843949ee627641c27b30}
\newcolumntype{Y}{>{\raggedright\arraybackslash}X}
\hypersetup{pdftitle={Tail-Spans and Differential Transcendence in Surreal and Surcomplex Hahn Fields},pdfauthor={Research article prepared for Vladimir Reshetnikov},pdfsubject={Algebraic relations, coefficient Galois actions, and explicit differential independence}}
\begin{document}
\hypersetup{pageanchor=false}
\begin{titlepage}
\centering
\vspace*{8mm}
{\large\scshape Research article for the Surreal project\par}
\vspace{11mm}
{\Huge\bfseries\color{Ink} Tail-Spans and\par Differential Transcendence\par}
\vspace{5mm}
{\LARGE in Surreal and Surcomplex\par Hahn Fields\par}
\vspace{8mm}
{\large September 22, 2026\par}
\vspace{8mm}
\begin{minipage}{.95\textwidth}
\begin{abstract}
We give an exact algebraic classification of finite tuples of Hahn series
whose coefficients use independent square roots. The decisive invariant is
not the full span of the coefficient vectors but their \emph{cofinite span}:
the subspace that remains after every finite deletion. If its dimension is
$d$ and $E$ is the finite set of exceptional vectors, the generated algebra
is exactly a polynomial algebra in $d$ variables over the multiquadratic
extension belonging to $E$. Consequently, both the transcendence degree and
all geometric components of the relation ideal are explicit.

A finite-difference argument converts independent coefficient sign changes
into vanishing conditions on the highest homogeneous part of a hypothetical
polynomial relation. This yields a continuum-sized, explicitly indexed
family of differentially algebraically independent \emph{actual surreal
numbers}, all with real-algebraic coefficients and positive integer
valuation support. In particular, the Berarducci--Mantova differential
transcendence degree of $\mathbb A_{\R}((t^{\Q}))$ over $\Q((t^{\Q}))$ is
$2^{\aleph_0}$, with $t=\omega^{-1}$.

For surcomplex analysis, the same mechanism proves that every mixed jet of
$\sum_{n\in A}t^n\sqrt{1-z/2^n}$ is algebraically independent over
$\C(z)((t^{\Q}))$, for every infinite $A\subseteq\N_{>0}$, with respect to
$d/dz$ and $t\,d/dt$. A continuum-sized family enjoys joint mixed-jet
independence on one common disk. We also give exact algebraic-approximation
criteria, finite witnesses, and limitations on coefficient-stream decision
procedures.
\end{abstract}
\end{minipage}
\vfill
\begin{minipage}{.95\textwidth}\small
\textbf{Research status.} Complete mathematical proofs are supplied relative
to standard Hahn-field and surreal normal-form foundations. The tail-span
classification and the explicit differential-independence constructions are
proposed contributions; a targeted literature search did not identify these
formulations, but priority is not certified. No named historical conjecture
is claimed solved. The article is not independently refereed. The Lean
ledger records proofs of cofinite-span stabilization and multilinear
detection (Lemmas~\ref{tail:lem:stabilization} and~\ref{tail:lem:multilinear});
the classification and differential-independence theorems remain pending.
Accompanying exact computations test finite identities and examples,
not the infinite theorems.
\end{minipage}
\end{titlepage}
\hypersetup{pageanchor=true}
\pagenumbering{roman}
\tableofcontents
\clearpage
\pagenumbering{arabic}

\section{The research question and the contributions}\label{tail:sec:intro}

An infinite Hahn sum can be far more algebraically complicated than any of
its coefficients. This is especially relevant to surreal mathematics:
a Conway normal form may use only real-algebraic coefficients and very
simple exponents, while the resulting number is not controlled by any
finite algebraic or differential equation over a substantial Hahn subfield.
For analytic coefficient functions the discrepancy is stronger still:
every coefficient may satisfy a quadratic equation on the same disk, while
the whole function and all of its mixed derivatives satisfy no polynomial
relation over a large rational-coefficient Hahn field.

The aim here is an exact structural result, not merely the production of
one transcendental series. We ask how all polynomial relations in a finite
family of such series can be read from the coefficient vectors. The answer
is a finite-dimensional invariant with a finite algebraic exceptional part.
This produces explicit differential independence, including inside the
actual surreal differential field.

\subsection{The main statements in compact form}
Let $M$ be a field of characteristic zero, let $\Gamma$ be an ordered
abelian group, and let $a_i\in M^\times$ have independent square classes.
Choose $r_i^2=a_i$ and distinct exponents $\gamma_i$ whose set is well
ordered. For $v_i\in M^r$, put
\begin{equation}\label{tail:eq:introY}
 Y=\sum_{i\in I} t^{\gamma_i}r_i v_i,
 \qquad B=M((t^\Gamma)),
 \qquad W=\bigcap_{F\subseteq I\text{ finite}}
       \Span_M\{v_i:i\notin F\}.
\end{equation}
Theorem~\ref{tail:thm:tail} proves that $E=\{i:v_i\notin W\}$ is finite and
\begin{equation}\label{tail:eq:intromain}
 B[Y_1,\ldots,Y_r]
   =M(r_i:i\in E)((t^\Gamma))[s_1,\ldots,s_d],
 \qquad d=\dim_M W,
\end{equation}
where the $s_j$ are algebraically independent over the displayed
coefficient field. The equality is an equality of subrings in a specified
ambient Hahn field, not merely an abstract comparison of dimensions.
Theorem~\ref{tail:thm:planes} then identifies the full relation ideal after its
finite splitting extension.

For actual surreal numbers, write
\[
 B_0=\Q((t^\Q)),\qquad H_0=\AR((t^\Q)),\qquad t=\omega^{-1},
\]
where $\AR$ is the field of real-algebraic numbers. We construct a family
$(\xi_\alpha)_{\alpha\in\{0,1\}^{\N}}\subset H_0$ such that
\begin{equation}\label{tail:eq:introBM}
 \bigl\{\BM^k\xi_\alpha:
        \alpha\in\{0,1\}^{\N},\ k\geq0\bigr\}
 \quad\text{is algebraically independent over }B_0.
\end{equation}
Thus $\dtrdeg_{B_0}H_0=\cf$, where $\cf=2^{\aleph_0}$
(Theorem~\ref{tail:thm:surreal}). This is a differential result about specified
set-sized subfields of $\No$, not an appeal to the size of the whole
proper class.

The analytic counterpart starts with the single explicit function
\begin{equation}\label{tail:eq:introF}
 F(z,t)=\sum_{n\geq1}t^n\sqrt{1-z/2^n}.
\end{equation}
Every square root is the branch equal to $1$ at $z=0$. On one common disk,
all mixed jets $\partial_z^j\EDer^kF$, with $\EDer=t\,d/dt$, are
algebraically independent over $\C(z)((t^\Q))$
(Theorem~\ref{tail:thm:mixed}). Theorem~\ref{tail:thm:continuumanalytic} supplies
continuum many jointly independent such functions.

\subsection{What was checked in the repository}
The repository was read at revision
\begin{center}\small\texttt{\repocommit}.\end{center}
The main README and documentation catalogue were inspected, together with
the opening spectral-theory article and the Laurent-birthday report
inventory~\cite{Repo,RepoMap}. The catalogue distinguishes common-domain
analytic Hahn functions from coefficientwise germs and separates
coefficientwise summation from the fine topology of $\No[i]$. Those
conventions are retained here.

The inspected catalogue already organizes substantial material on
polynomial algebra, analytic geometry, residues, differential equations,
rank-one geometry, and finite spectral theory. The present article instead
centers on \emph{algebraic relations among infinitely many algebraic
coefficients}, and on the differential transcendence they generate. This
is a targeted coverage audit, not a claim to have read every archived
manuscript or every Lean declaration. In particular, an unfinished Lean
bridge is not treated as an open mathematical theorem.

\subsection{Attribution and novelty boundaries}
Hahn support calculus, finite separable base change, and multiquadratic
Galois theory are background, not claimed innovations. We give proofs of
the algebraic ingredients to make the new argument independently auditable;
standard Galois and Kummer references are the Stacks Project
\cite{StacksGalois,StacksKummer}. Generalized-series operator methods have
an extensive literature, including van der Hoeven~\cite{vdH}. Conway
normal forms and their use as Hahn fields are reviewed, for example, by
Kaplan--Krapp--Serra~\cite{KKS}. The existence, strong additivity, real
constant field, and normalization $\BM\omega=1$ of the derivation used
here are due to Berarducci--Mantova~\cite{BM}.

The proposed contribution is the \emph{combined exact classification} in
\eqref{tail:eq:intromain}, its elementary polarization proof, and the explicit
ordinary and mixed differential-independence constructions. Differential
transcendence itself is a classical subject; methods based on differential
fields and functional equations, for example, appear in
Mijajlovi\'c--Male\v sevi\'c~\cite{MM}. No claim is made that the general
idea of producing transcendence from growing coefficient extensions is
new. The searches were targeted rather than exhaustive, and neither a
priority theorem nor a resolution of a previously named open problem is
asserted.

\subsection{Two places where this report meets its companions}
The restriction of the Berarducci--Mantova derivation to a
rational-exponent Hahn workspace is \emph{not} proved here. The companion
report on differential algebra and differential equations~\cite{RepoDiffEq}
fixes that workspace and proves that formula, and
Section~\ref{tail:sec:surreal} cites it. Only the further restriction from
real to real-algebraic coefficients, which is immediate, is checked here.

In the other direction, the companion report on dynamics and normal
forms~\cite{RepoDyn} carries a warning, \texttt{dyn:warn:diffalg}, that its
own transcendental solutions must \emph{not} be called differentially
transcendental, since they satisfy a first-order linear equation over its
base $\C(w)((t^\Gamma))$. Theorem~\ref{tail:thm:mixed} obtains the
differential conclusion that warning declines, for different functions,
over the $\Gamma=\Q$ case of a base of the same shape and on one fixed
disk with no shrinking. Remark~\ref{tail:rem:dynupgrade} states the
relation precisely and does not weaken the warning, which remains true of
the functions it is about.

\section{Hahn fields, coefficient extensions, and summation}\label{tail:sec:hahn}

\subsection{Set-sized conventions}
All fields, groups, index families, and vector spaces in the algebraic
proofs are sets. A Hahn series over a field $M$ with value group $\Gamma$
is a formal expression
\[
 f=\sum_{\gamma\in S}c_\gamma t^\gamma,
 \qquad c_\gamma\in M,\quad S\subseteq\Gamma\text{ well ordered}.
\]
Its support consists of the exponents with nonzero coefficient. Addition
is coefficientwise, multiplication is convolution, and
$v(f)=\min\supp(f)$ for $f\ne0$, with $v(0)=+\infty$.
The resulting field is denoted $M((t^\Gamma))$.

A family is \emph{Hahn summable} when its supports have well-ordered union
and only finitely many family members contribute at each exponent. This
is a support condition. It does not mean convergence of a sequence of
partial sums in the fine topology of the entire surreal field.
We use the standard Hahn support facts: finite sums of well-ordered
subsets are well ordered, each fixed exponent has only finitely many
representations from a pair of such supports, and the monoid generated
by a well-ordered positive set has the corresponding finite-word
property. These are the usual Neumann support lemmas; their operator
uses are treated in~\cite{vdH}.

For $\Gamma=\Q$, the map
\begin{equation}\label{tail:eq:embedding}
 \sum_q c_qt^q\longmapsto\sum_q c_q\omega^{-q}
\end{equation}
embeds $\R((t^\Q))$ into $\No$ and, coefficientwise in real and imaginary
parts, $\C((t^\Q))$ into $\No[i]$. This is an instance of the standard
Conway normal-form correspondence~\cite{KKS}. No construction of the
entire surreal field, or of its real closedness, is part of this paper.

\subsection{Finite coefficient extensions commute with Hahn series}
\begin{lemma}[Finite base change]\label{tail:lem:basechange}
Let $F/M$ be a finite field extension of degree $e$, with basis
$u_1,\ldots,u_e$. Inside a common coefficient extension,
\[
 F((t^\Gamma))=M((t^\Gamma))F,
 \qquad [F((t^\Gamma)):M((t^\Gamma))]=e.
\]
The same $u_1,\ldots,u_e$ form a basis after Hahn base change. If $F/M$
is finite Galois, its automorphisms act coefficientwise and give the full
Galois group of the base-changed extension.
\end{lemma}
\begin{proof}
Expand each coefficient of $f\in F((t^\Gamma))$ in the $M$-basis:
$c_\gamma=\sum_j m_{j,\gamma}u_j$. For each $j$, the support of
$\sum_\gamma m_{j,\gamma}t^\gamma$ is contained in $\supp(f)$.
Thus $f=\sum_j u_j f_j$ with $f_j\in M((t^\Gamma))$.
Conversely, a finite sum of this kind is an $F$-coefficient Hahn series.
If $\sum_j u_j f_j=0$, coefficient comparison and the original basis
independence imply that every $f_j$ is zero. This proves the degree and
basis assertions. Each coefficient automorphism preserves supports and
convolution, and the $e$ distinct coefficient automorphisms exhaust the
Galois group when $F/M$ is Galois.
\end{proof}

\begin{warning}\label{tail:warn:infinitebase}
For an infinite algebraic extension $L/M$, the analogous equality
$L((t^\Gamma))=M((t^\Gamma))L$ generally fails. The field on the right is
the union of the finite constant extensions from Lemma~\ref{tail:lem:basechange};
a Hahn series on the left can involve coefficients generating an infinite
extension. Section~\ref{tail:sec:coefficient} makes the obstruction exact, and
Section~\ref{tail:sec:surreal} shows that the failure can have maximal
differential transcendence degree in a countable-coefficient workspace.
\end{warning}

\subsection{Independent square classes and independent sign changes}
\begin{definition}
A family $(a_i)_{i\in I}\subset M^\times$ has \emph{independent square
classes} if, for every nonempty finite $J\subseteq I$,
$\prod_{i\in J}a_i$ is not a square in $M$.
Fix square roots $r_i$ in an algebraic closure and put
$L=M(r_i:i\in I)$.
\end{definition}

\begin{lemma}[Independent sign changes]\label{tail:lem:signs}
Assume $\operatorname{char}M\ne2$. For finite $J\subseteq I$, the extension
$M(r_i:i\in J)/M$ has degree
$2^{|J|}$ and basis $\{\prod_{i\in U}r_i:U\subseteq J\}$.
Every assignment $r_i\mapsto\varepsilon_i r_i$, $\varepsilon_i\in\{1,-1\}$,
extends to an $M$-automorphism of $L$. It extends coefficientwise to an
automorphism of $L((t^\Gamma))$ fixing $M((t^\Gamma))$.
\end{lemma}
\begin{proof}
We give the usual elementary multiquadratic argument. Suppose the claim
holds for $r_1,\ldots,r_m$, and write $E_m=M(r_1,\ldots,r_m)$.
If $u^2=a_{m+1}$ for some $u\in E_m$, each sign automorphism $\sigma$
satisfies $\sigma(u)=\pm u$. Hence $u$ is a simultaneous eigenvector for
the sign group. The basis monomials $r_U$ have distinct sign characters,
so their simultaneous eigenspaces are exactly the one-dimensional spaces
$Mr_U$. Thus $u=cr_U$, contradicting independence of the square classes.
Adjoining $r_{m+1}$ therefore doubles the degree. The old sign changes
extend by fixing $r_{m+1}$, and the new sign change fixes $E_m$.
This proves the finite assertion by induction, starting with $m=0$.

Every element of $L$ belongs to a finite subextension. The automorphisms
given by one fixed sign assignment are compatible on all such
subextensions and define an automorphism of their union. Coefficientwise
extension to Hahn series preserves zero coefficients, support,
addition, and the finite convolution sums, proving the last assertion.
\end{proof}

The characteristic restriction is necessary. In $M=\mathbb F_2(u,v)$,
the classes of $u,v,u+v$ are independent: each nonempty product has odd
order at one of these distinct prime divisors. But
$\sqrt{u+v}=\sqrt u+\sqrt v$, so adjoining the three roots gives degree
at most $4$, not $2^3$.

\begin{remark}[No order-preserving automorphism is asserted]\label{tail:rem:noorder}
When $M=\Q$ and the $r_i$ are positive real square roots, a sign change
need not preserve the order of the multiquadratic subfield. We use it
only as a field automorphism of that subfield and its Hahn field. We do
\emph{not} claim that it extends to an automorphism of the ordered field
of all real-algebraic numbers, or to an automorphism of $\No$. The proofs
below need neither extension.
\end{remark}

\section{The exact coefficient-field criterion}\label{tail:sec:coefficient}

The following preliminary is a useful classical Galois-orbit argument. It
isolates the scalar algebraicity issue before the genuinely multivariable
classification.

\begin{theorem}[Coefficient-field criterion]\label{tail:thm:coefficient}
Let $M$ have characteristic zero, let $\overline M$ be an algebraic
closure, and let
\[
 f=\sum_\gamma c_\gamma t^\gamma\in\overline M((t^\Gamma)),
 \qquad F=M(c_\gamma:\gamma\in\supp(f)),
 \qquad B=M((t^\Gamma)).
\]
Then $f$ is algebraic over $B$ if and only if $F/M$ is finite. In that case
\begin{equation}\label{tail:eq:coeffdegree}
 B(f)=F((t^\Gamma)),\qquad [B(f):B]=[F:M].
\end{equation}
Consequently, the relative algebraic closure of $B$ inside
$\overline M((t^\Gamma))$ is
\begin{equation}\label{tail:eq:relativeunion}
 \mathcal A=\bigcup_{F/M\text{ finite},\ F\subseteq\overline M}
                  F((t^\Gamma)).
\end{equation}
\end{theorem}
\begin{proof}
Assume first that $F/M$ is finite. Lemma~\ref{tail:lem:basechange} gives
$[B(f):B]\leq[F:M]$. Every $M$-embedding of $F$ into $\overline M$ extends
to an automorphism of $\overline M$; this is the standard extension
property for normal algebraic extensions
\cite[Lemma~9.15.7(2)]{StacksNormal}, applied to $\overline M/M$.
Characteristic zero makes $F/M$ separable, so there are exactly $[F:M]$
such embeddings. Distinct embeddings differ on a coefficient of $f$,
since those coefficients generate $F$. Their coefficientwise extensions
therefore give $[F:M]$ distinct conjugates of $f$ over $B$. The reverse
degree inequality follows. Since $B(f)\subseteq F((t^\Gamma))$ and both
extensions have the same finite degree, equality holds.

Conversely, suppose $f$ has degree $D$ over $B$. Choose any finite
subextension $F_0/M$ generated by finitely many coefficients of $f$.
Its distinct embeddings into $\overline M$ extend as above and produce
at least $[F_0:M]$ distinct conjugates of $f$. Hence $[F_0:M]\leq D$.
The degrees of all such $F_0$ form a bounded nonempty set of positive
integers. Choose one, $F_*$, of maximal degree. Adjoining any coefficient
$c_\gamma$ gives another finite coefficient-generated field and cannot
increase that degree. Therefore $c_\gamma\in F_*$ for every $\gamma$,
so $F=F_*$ is finite. Formula~\eqref{tail:eq:relativeunion} follows immediately.
\end{proof}

\begin{example}
Let $p_n$ be the $n$th prime and choose the positive roots. Then
\[
 x=\sum_{n\geq1}\sqrt{p_n}\,t^n
\]
is transcendental over $\Q((t^\Q))$: its coefficients generate an
infinite multiquadratic extension. Every finite truncation is algebraic,
and its first $N$ terms generate a degree-$2^N$ constant extension after
Hahn base change. Section~\ref{tail:sec:surreal} proves the substantially
stronger assertion that $x$ is differentially transcendental and belongs
to large explicit families of differentially independent surreals.
\end{example}

\section{Cofinite spans and the polarization mechanism}\label{tail:sec:mechanism}

\subsection{The span that survives every finite deletion}
Let $(v_i)_{i\in I}$ be any family in a finite-dimensional $M$-space $V$.
No finite alphabet, boundedness, periodicity, or computability hypothesis
is imposed on the $v_i$.

\begin{definition}[Cofinite span]
The \emph{cofinite span} of the family is
\[
 W=\bigcap_{F\subseteq I\text{ finite}}
             \Span_M\{v_i:i\notin F\}.
\]
The term ``tail-span'' will also be used. Here ``tail'' means cofinite:
for an arbitrary ordered support it does not mean deletion of an arbitrary
transfinite initial segment.
\end{definition}

\begin{lemma}[Stabilization and exceptions]\label{tail:lem:stabilization}
There is a finite $F_0\subseteq I$ such that
$W=\Span_M\{v_i:i\notin F_0\}$. The set
\[
 E=\{i\in I:v_i\notin W\}
\]
is finite. Moreover, every cofinite subfamily of $(v_i)_{i\in I\setminus E}$
spans $W$.
\end{lemma}
\begin{proof}
Among the dimensions of spans obtained by finite deletion, choose a
minimum, attained by $F_0$. Deleting any additional finite set from that
family gives a subspace of the same dimension, hence the same subspace.
For any finite $F$, the span obtained by deleting $F\cup F_0$ is therefore
the $F_0$-span and is contained in the $F$-span. This proves that the
intersection equals the $F_0$-span. In particular every $i\notin F_0$ has
$v_i\in W$, so $E\subseteq F_0$. Finally, deleting $E$ and any other finite
set leaves a family contained in $W$ whose span contains the defining
intersection $W$. Its span is exactly $W$.
\end{proof}

\subsection{Multilinear forms detected by distinct indices}
\begin{lemma}[Cofinite spanning detects multilinear forms]\label{tail:lem:multilinear}
Suppose every cofinite subfamily of $(w_i)_{i\in J}$ spans a finite-dimensional
space $W$. Let $K/M$ be a field extension, and let $T$ be a $D$-linear form
on $W\otimes_M K$, with $D\geq1$. If
\[
 T(w_{i_1},\ldots,w_{i_D})=0
 \quad\text{whenever }i_1,\ldots,i_D\text{ are pairwise distinct},
\]
then $T=0$.
\end{lemma}
\begin{proof}
If $W=0$, there is nothing to prove. Otherwise $J$ is infinite, since a
finite family would leave the zero span after a finite deletion.
Fix distinct $i_1,\ldots,i_{D-1}$. The remaining vectors span $W$, and
therefore span its scalar extension over $K$. Linearity frees the last
argument: the expression vanishes for any value of that argument.
Now fix only $D-2$ distinct indices and an arbitrary last argument.
The possible next-to-last indexed vectors outside these indices still
span $W$, so that argument is also free. Repeating this procedure frees
all $D$ arguments. For $D=1$ the first spanning step already finishes the
proof.
\end{proof}

\begin{lemma}[Top-degree mixed difference]\label{tail:lem:difference}
Let $K$ have characteristic zero and let the nonzero polynomial
$P\in K[X_1,\ldots,X_d]$ have total degree $D$, with homogeneous part $P_D$.
For vectors $u_1,\ldots,u_D\in K^d$, define
$\Delta_uP(X)=P(X+u)-P(X)$. Then
\begin{equation}\label{tail:eq:diffidentity}
 \Delta_{u_1}\cdots\Delta_{u_D}P(X)
  =D_{u_1}\cdots D_{u_D}P_D
  =D!\,T(u_1,\ldots,u_D),
\end{equation}
where $D_u$ is directional differentiation and $T$ is the symmetric
polarization satisfying $P_D(X)=T(X,\ldots,X)$. The result is independent
of $X$. For $D=0$, the empty operator products are the identity and
$T=P_0$ is a scalar.
\end{lemma}
\begin{proof}
The case $D=0$ is immediate. For $D\ge1$, the finite polynomial Taylor
identity gives, on polynomials of degree at most $D$,
\[
 \Delta_u=\sum_{m=1}^D\frac{D_u^m}{m!}.
\]
In a product of $D$ such differences, every term has total differential
order at least $D$. A term of order greater than $D$ annihilates $P$.
The only term of order exactly $D$ chooses the first derivative in every
factor; it annihilates every homogeneous part below degree $D$.
This proves the first equality, even if some directions are zero or a
directional derivative of $P_D$ vanishes. Expanding
$P_D(\lambda_1u_1+\cdots+\lambda_Du_D)$ and extracting the coefficient of
$\lambda_1\cdots\lambda_D$ gives the second. Characteristic zero allows
division by $D!$, and identifies a homogeneous polynomial uniquely with
its symmetric polarization.
\end{proof}

\subsection{Why coefficient sign changes defeat polynomial relations}
\begin{proposition}[Sign-orbit independence]\label{tail:prop:orbit}
Let $M$ have characteristic zero, let $B=M((t^\Gamma))$, let
$(a_i)_{i\in I}$ have independent square classes, and let $F\subset I$
be finite. Set
$B_F=M(r_i:i\in F)((t^\Gamma))$. Suppose $w_i\in M^d$ for $i\notin F$
and every cofinite subfamily of these vectors spans $M^d$. If the
$\gamma_i$ are distinct with well-ordered image, then the coordinates of
\[
 s=\sum_{i\notin F}t^{\gamma_i}r_iw_i
\]
are algebraically independent over $B_F$.
\end{proposition}
\begin{proof}
For $d=0$ the assertion is the empty independence assertion. Otherwise,
suppose $P(s)=0$ for a nonzero polynomial $P\in B_F[X_1,\ldots,X_d]$ of
positive degree $D$. A nonzero constant cannot vanish, so all possible
relations have this form. Put $b_i=t^{\gamma_i}r_i$.

For each $i\notin F$, the coefficient sign change $\sigma_i$ fixes $B_F$
and sends $s$ to $s-2b_iw_i$. Distinct such sign changes commute and
leave one another's $b_i$ unchanged. Applying their products to $P(s)=0$
shows that
\[
 P\left(s-2\sum_{j\in J}b_{i_j}w_{i_j}\right)=0
 \quad\text{for every }J\subseteq\{1,\ldots,D\},
\]
whenever $i_1,\ldots,i_D$ are distinct and outside $F$.
Take the alternating sum of these $2^D$ identities. By
Lemma~\ref{tail:lem:difference}, it is
\begin{equation}\label{tail:eq:polarzero}
 0=(-2)^D D!\left(\prod_{j=1}^D b_{i_j}\right)
                  T(w_{i_1},\ldots,w_{i_D}),
\end{equation}
where $T$ is the polarization of $P_D$, over the coefficient field $B_F$.
All the scalar factors preceding $T$ are nonzero. Thus $T$ vanishes on
all distinct indexed tuples. Lemma~\ref{tail:lem:multilinear} gives $T=0$,
contradicting $P_D\ne0$. No infinite sequence of Galois operations, and
no limiting argument, is used: a relation of degree $D$ is destroyed by
finitely many sign changes.
\end{proof}

\section{The tail-span classification}\label{tail:sec:classification}

\subsection{The generated algebra and its algebraic part}
\begin{theorem}[Exact tail-span classification]\label{tail:thm:tail}
Let $M$ have characteristic zero and let $\Gamma$ be a set-sized ordered
abelian group. Let $(a_i)_{i\in I}\subset M^\times$ have independent square
classes, choose $r_i^2=a_i$, and put $L=M(r_i:i\in I)$.
Let $\gamma_i\in\Gamma$ be distinct with well-ordered image and let
$v_i\in M^r$. In $H=L((t^\Gamma))$, define
\[
 Y=\sum_i t^{\gamma_i}r_iv_i,\qquad B=M((t^\Gamma)),
 \qquad W=\bigcap_{F\subseteq I\text{ finite}}
                     \Span_M\{v_i:i\notin F\}.
\]
Set $d=\dim_M W$, $E=\{i:v_i\notin W\}$, and
\[
 p=\sum_{i\in E}t^{\gamma_i}r_iv_i,
 \qquad B_E=M(r_i:i\in E)((t^\Gamma)).
\]
Then $E$ is finite. For any $M$-basis of $W$, let $s_1,\ldots,s_d$ be
the coordinates of $Y-p$. These coordinates satisfy
\begin{equation}\label{tail:eq:algebraclassification}
 B[Y_1,\ldots,Y_r]=B_E[s_1,\ldots,s_d],
\end{equation}
and are algebraically independent over $B_E$. In particular,
\begin{align}
 B(Y_1,\ldots,Y_r)&=B_E(s_1,\ldots,s_d),\label{tail:eq:fieldclassification}\\
 \trdeg_B B(Y_1,\ldots,Y_r)&=d,\label{tail:eq:trdegree}\\
 [B_E:B]&=2^{|E|}.\label{tail:eq:exceptiondegree}
\end{align}
The relative algebraic closure of $B$ in $B(Y_1,\ldots,Y_r)$ is exactly
$B_E$.
\end{theorem}
\begin{proof}
By Lemma~\ref{tail:lem:stabilization}, $E$ is finite and the vectors outside
$E$ lie in $W$, with every cofinite subfamily spanning $W$. In coordinates
on $W$, Proposition~\ref{tail:prop:orbit} proves algebraic independence of
$s_1,\ldots,s_d$ over $B_E$. It remains to show that all of $B_E$ is
already present in $B[Y_1,\ldots,Y_r]$; this is stronger than the
transcendence-degree assertion alone.

If $E$ is empty, $B_E=B$ and the assertion is immediate. Otherwise pass
to the quotient $M^r/W$. Each $v_i$ with $i\in E$ has nonzero image.
There is an $M$-linear functional $\ell:M^r\to M$ vanishing on $W$ and
nonzero on every such $v_i$. Indeed, a finite union of proper hyperplanes
cannot exhaust the dual of a finite-dimensional vector space over the
infinite field $M$. To see this directly, choose coordinates on the quotient
dual. Each nonzero image of $v_i$ defines a nonzero linear polynomial
there, and their product is nonzero. A nonzero polynomial over an
infinite field cannot vanish at every point, by induction on its number
of variables. A point where this product is nonzero supplies $\ell$.
Thus
\begin{equation}\label{tail:eq:primitiveeta}
 \eta=\ell(Y)=\sum_{i\in E}\ell(v_i)t^{\gamma_i}r_i
             \in B[Y_1,\ldots,Y_r]\cap B_E.
\end{equation}
The extension $B_E/B$ is multiquadratic of degree $2^{|E|}$ by
Lemmas~\ref{tail:lem:basechange} and~\ref{tail:lem:signs}. Every sign assignment
on $E$ gives a different conjugate of $\eta$: at an index where two
assignments differ, the coefficient of its distinct exponent $\gamma_i$
is nonzero and changes sign. Hence $\eta$ has $2^{|E|}$ distinct
conjugates, and $B(\eta)=B_E$.

Because $\eta$ is algebraic over the field $B$, its polynomial algebra
$B[\eta]$ is already the field $B(\eta)$. Consequently
$B_E\subseteq B[Y_1,\ldots,Y_r]$. All coordinates of $p$ therefore
belong to this generated algebra. Extend the chosen basis of $W$ to a
basis of $M^r$; its dual coordinate forms recover the $s_j$ from $Y-p$.
Thus these coordinates also belong to the generated algebra. This proves the
inclusion from right
to left in~\eqref{tail:eq:algebraclassification}. The reverse inclusion follows
from $Y=p+Us$, where the columns of the matrix $U$ are the chosen basis
vectors of $W$. The field and degree assertions follow.

Finally, a purely transcendental extension $K(s_1,\ldots,s_d)$ contains
no elements algebraic over $K$ except $K$ itself. For completeness, if a
nonconstant rational function depends on some $s_j$, regard it as a
nonconstant rational function in that variable over the field of the
others; the variable is then algebraic over that rational function,
which proves the rational function is transcendental over that field.
An element algebraic over $B$ is also algebraic over $B_E$, so it must
belong to $B_E$. Conversely every element of $B_E$ is algebraic over $B$.
\end{proof}

\subsection{All geometric components of the relation ideal}
\begin{theorem}[Relation ideal as disjoint affine planes]\label{tail:thm:planes}
In Theorem~\ref{tail:thm:tail}, let
\[
 I_Y=\ker\bigl(B[Z_1,\ldots,Z_r]\longrightarrow H,
                     \ Z_j\longmapsto Y_j\bigr).
\]
Choose an $M$-linear surjection $q:M^r\to M^{r-d}$ with kernel $W$.
For $\varepsilon\in\{1,-1\}^E$, put
\[
 p_\varepsilon=\sum_{i\in E}
                 \varepsilon_i t^{\gamma_i}r_iv_i.
\]
Then
\begin{equation}\label{tail:eq:planes}
 I_YB_E[Z_1,\ldots,Z_r]
   =\bigcap_{\varepsilon\in\{1,-1\}^E}
                 \bigl(q(Z)-q(p_\varepsilon)\bigr).
\end{equation}
Each ideal on the right is generated by $r-d$ independent affine linear
forms. The $2^{|E|}$ ideals are pairwise comaximal. Thus the base-changed
coordinate algebra is the product of $2^{|E|}$ polynomial algebras in
$d$ variables over $B_E$.
\end{theorem}
\begin{proof}
Theorem~\ref{tail:thm:tail} identifies the coordinate algebra over $B$ with
$B_E[s_1,\ldots,s_d]$. Finite Galois base change gives
\[
 B_E\otimes_B B_E\cong
       \prod_{\sigma\in\Gal(B_E/B)}B_E,
 \qquad a\otimes b\longmapsto(a\sigma(b))_\sigma.
\]
For $E\ne\varnothing$, this splitting follows directly from the primitive
element $\eta$ in \eqref{tail:eq:primitiveeta}: its minimal polynomial over
$B$ splits over $B_E$ into distinct factors $T-\sigma(\eta)$, and the
Chinese remainder theorem identifies $B_E[T]/(m_\eta)$ with the displayed
product. For $E=\varnothing$, the identity is simply $B\otimes_B B\cong B$.
Adjoining the independent $s_j$ and using $Y=p+Us$, the component indexed
by $\sigma$ sends $Z$ to $\sigma(p)+Us$. Its kernel is precisely
$(q(Z)-q(\sigma(p)))$. The sign assignments index all $\sigma$,
proving~\eqref{tail:eq:planes}.

If $\varepsilon\ne\varepsilon'$, then
$q(p_\varepsilon)-q(p_{\varepsilon'})\ne0$. To see this, choose an index
where the signs differ. At its distinct exponent the coefficient vector
is a nonzero multiple of $q(v_i)$, and $q(v_i)\ne0$ by the definition of
$E$. At least one quotient coordinate therefore gives two distinct
constants. Subtracting the corresponding affine generators gives a
nonzero element of the coefficient field $B_E$, hence a unit. The ideals
are comaximal. The Chinese remainder theorem gives the stated product;
in particular the components are disjoint and reduced.
\end{proof}

\begin{remark}
Over $B$, the ideal $I_Y$ is prime, because it is the kernel of a map to
a field. Its decomposition after $B_E$-base change is not a contradiction:
the finite algebraic constant field splits into its conjugate factors.
The cofinite span determines the affine dimension, while the exceptional
vectors determine the entire finite splitting degree.
\end{remark}

\section{Examples, exact relations, and finite witnesses}\label{tail:sec:examples}

\subsection{One exceptional coefficient produces a quadratic relation}
Take $I=\N_{>0}$, $\gamma_n=n$, $v_1=(1,0)$, and $v_n=(1,1)$ for
$n\geq2$. Put $b_n=t^nr_n$ and $S=\sum_{n\geq2}b_n$. Then
\[
 Y_1=b_1+S,\qquad Y_2=S,\qquad
 W=M(1,1),\quad E=\{1\}.
\]
The classification becomes
\[
 B[Y_1,Y_2]=B(r_1)[S],\qquad \trdeg_B B(Y_1,Y_2)=1.
\]
The complete relation ideal is
\begin{equation}\label{tail:eq:quadraticexample}
 I_Y=\bigl((Z_1-Z_2)^2-t^2a_1\bigr).
\end{equation}
Indeed, using $U=Z_1-Z_2$ and $V=Z_2$, the quotient is
$B[U]/(U^2-t^2a_1)[V]$, and the quadratic is irreducible because $r_1$
remains quadratic after finite base change. Over $B(r_1)$ it splits into
the two parallel lines $Z_1-Z_2=\pm tr_1$.

This example explains why simply declaring the tuple ``independent up to
finitely many terms'' would lose important information. The finite term
is precisely what creates the nontrivial algebraic constant field.

\subsection{Indicator families and recurrence patterns}
For subsets $A_1,\ldots,A_r\subseteq\N_{>0}$, let
\[
 Y_j=\sum_{n\in A_j}t^nr_n,\qquad
 v_n=(\mathbf1_{A_1}(n),\ldots,\mathbf1_{A_r}(n)).
\]
Only finitely many vector patterns are possible. The cofinite span is
exactly the span over $M$ of the patterns occurring infinitely often.
Patterns occurring only finitely often can all be removed, whereas an
infinitely recurrent pattern survives every finite deletion. Thus
\[
 \trdeg_B B(Y_1,\ldots,Y_r)
   =\dim_M\Span\{\text{infinitely recurrent patterns}\}.
\]
Equivalently, this is the rank of the indicator sequences modulo the
subspace of finitely supported sequences. For example, pairwise almost
disjoint infinite sets give full rank, because each coordinate unit vector
occurs infinitely often outside the finite overlaps.

\subsection{Vandermonde vectors need no recurring finite alphabet}
For $v_n=(1,n,\ldots,n^{r-1})$ and distinct positive integers $n$, every
$r$ such vectors are independent. Their determinant is the nonzero
Vandermonde product $\prod_{u<v}(n_v-n_u)$. Therefore every infinite
subfamily has cofinite span $M^r$. Theorem~\ref{tail:thm:tail} gives
\[
 \sum_n t^nr_n,\quad
 \sum_n nt^nr_n,\quad\ldots,\quad
 \sum_n n^{r-1}t^nr_n
\]
as an algebraically independent tuple over $M((t^\Gamma))$ whenever the
positive integer exponents are available and the chosen index set is
infinite. These are exactly the successive Euler derivatives used below.

\subsection{Finite witnesses against a bounded-degree relation}
\begin{corollary}\label{tail:cor:finitewitness}
In Proposition~\ref{tail:prop:orbit}, suppose $d>0$. For any degree $D\geq1$,
choose $D$ pairwise disjoint finite index sets, each supplying a basis of
$M^d$. Then a nonzero polynomial of degree $D$ cannot vanish at every
point in the sign orbit of $s$ generated by these indices. Some choice
of one index from each basis yields a nonzero $D$-fold mixed difference.
Thus at most $dD$ sign-change indices suffice for such a witness.
\end{corollary}
\begin{proof}
If the highest polarization vanished on every choice of one basis vector
from each of the $D$ bases, multilinearity would make it vanish on all of
$(M^d)^D$. Choose a tuple where it does not vanish. The indices are
pairwise distinct because the bases have disjoint index sets.
Equation~\eqref{tail:eq:diffidentity} then gives a nonzero alternating sum of
$2^D$ orbit evaluations. Such disjoint bases exist because every cofinite
subfamily spans $M^d$.
\end{proof}

This is a degree-sensitive finite certificate inside the proof. It is not
an algorithm deciding arbitrary infinite coefficient streams: locating
tail-spanning bases or certifying that they persist forever is additional
information. Section~\ref{tail:sec:effective} makes that distinction explicit.

\section{Actual surreal numbers of maximal differential independence}\label{tail:sec:surreal}

\subsection{A specified differential Hahn workspace}
Let $\AR=\overline\Q\cap\R$ be the countable field of real-algebraic
numbers. Via~\eqref{tail:eq:embedding}, regard
\begin{equation}\label{tail:eq:realworkspaces}
 B_0=\Q((t^\Q))\subset H_0=\AR((t^\Q))\subset\No,
 \qquad t=\omega^{-1},
\end{equation}
as concrete subfields of the surreal field. The Euler derivation is
\begin{equation}\label{tail:eq:euler}
 \EDer\left(\sum_q c_qt^q\right)=\sum_q q c_qt^q.
\end{equation}
It preserves supports up to deletion, so it is well defined and strongly
additive. The identity $q_1+q_2=q$ in each convolution coefficient proves
the Leibniz rule.

The normalized Berarducci--Mantova derivation $\BM$ on $\No$ has real
constant field, is strongly additive, and satisfies
$\BM\omega=1$~\cite{BM}. Its restriction to a rational-monomial Hahn
workspace is not reproved here. The companion report on differential
algebra and differential equations~\cite{RepoDiffEq} already fixes that
workspace: its section ``The rational-exponent workspace'' introduces
$\C((t^\Q))$ and $\R((t^\Q))$ with $t=\omega^{-1}$ at its equation
\texttt{diff:eq:HahnQ}, and its Proposition \texttt{diff:prop:hahnderiv},
at its equation \texttt{diff:eq:hahnderiv}, proves for every additive
subgroup $\Gamma\subseteq\R$ with $1\in\Gamma$ that $\C((t^\Gamma))$ is
$\BM$-stable, has constant field exactly $\C$, and satisfies
\[
 \BM\left(\sum_q c_qt^q\right)=-\sum_q q\,c_qt^{q+1}.
\]
Taking $\Gamma=\Q$ and comparing with~\eqref{tail:eq:euler} gives
\begin{equation}\label{tail:eq:BMrestriction}
 \BM(t^q)=-q t^{q+1},\qquad
 \BM\big|_{H_0}=-t\EDer.
\end{equation}
That proposition is stated for the complex coefficient field, which
contains ours; the further restriction to $H_0$ and $B_0$ needs only
that $-t\EDer$ multiplies each coefficient by a rational number and shifts
its exponent by $1$, so $\AR$ coefficients stay in $\AR$ and $\Q$
coefficients stay in $\Q$. Thus both $B_0$ and $H_0$ are differential
subfields. Their constant
fields are respectively $\Q$ and $\AR$: coefficient comparison in
$\EDer f=0$ leaves only exponent zero. The large differential
transcendence degree below is therefore not caused by adjoining
transcendental constants. This restriction
uses only rational powers; no unproved formula for the derivative of an
arbitrary Conway monomial is needed. The same companion report's warning
\texttt{diff:warn:monomialmap} records why the Conway monomial map
$\gamma\mapsto\omega^{-\gamma}$ must not be silently replaced by
$\exp(-\gamma\log\omega)$ at a general surreal exponent $\gamma$; both
readings agree at the ordinary rational exponents used here, and nowhere
below is a monomial with a non-real exponent differentiated.

Induction, or evaluation on every monomial followed by strong additivity,
gives the useful triangular identity
\begin{equation}\label{tail:eq:BMtriangular}
 \BM^k=(-1)^k t^k\EDer(\EDer+1)\cdots(\EDer+k-1),
 \qquad k\geq0.
\end{equation}
At $k=0$, the product is the identity. The leading coefficient in terms
of $\EDer^k$ is the nonzero base-field element $(-1)^kt^k$.
For every finite collection of elements and every finite jet order,
their $\BM$-jets and Euler jets are related by invertible triangular
linear transformations over $B_0$.

\subsection{One number and every infinite subsequence}
Let $p_n$ be the $n$th prime and $r_n=\sqrt{p_n}>0$.
Unique factorization in $\Z$ shows that the square classes of the $p_n$
are independent in $\Q^\times/\Q^{\times2}$: a nonempty finite product
has an odd exponent at every selected prime. For $A\subseteq\N_{>0}$,
define the actual surreal number
\begin{equation}\label{tail:eq:xiA}
 \xi_A=\sum_{n\in A}\sqrt{p_n}\,\omega^{-n}
      =\sum_{n\in A}r_nt^n.
\end{equation}
It is positive and infinitesimal when $A$ is nonempty.

\begin{theorem}[Every infinite subseries is differentially transcendental]
\label{tail:thm:onesurreal}
For every infinite $A\subseteq\N_{>0}$, the family
$\{\EDer^k\xi_A:k\geq0\}$ is algebraically independent over $B_0$.
The same is true of $\{\BM^k\xi_A:k\geq0\}$.
\end{theorem}
\begin{proof}
For any $K\geq0$,
\[
 (\xi_A,\EDer\xi_A,\ldots,\EDer^K\xi_A)
   =\sum_{n\in A}t^nr_n(1,n,\ldots,n^K).
\]
Every $K+1$ distinct coefficient vectors are a Vandermonde basis over
$\Q$. Consequently every cofinite subfamily spans $\Q^{K+1}$.
Apply Theorem~\ref{tail:thm:tail} with $W=\Q^{K+1}$ and $E=\varnothing$.
It gives independence of the finite jet tuple. Every polynomial in the
infinite jet family uses only finitely many jets, proving independence
of the whole family. Formula~\eqref{tail:eq:BMtriangular} transfers this
independence to the $\BM$-jets.
\end{proof}

\begin{remark}
The proof took place in the intermediate multiquadratic Hahn field
$\Q(\sqrt{p_n}:n\geq1)((t^\Q))$. Its coefficient sign automorphisms need
not extend to $H_0$ or $\No$; Remark~\ref{tail:rem:noorder} applies. Polynomial
independence proved in this intermediate field remains true for the same
elements in $H_0$ and $\No$, over the same base $B_0$.
\end{remark}

\subsection{An explicit continuum-sized almost disjoint family}
For a nonempty binary string $s=(s_1,\ldots,s_m)$, define
\begin{equation}\label{tail:eq:binarycode}
 N(s)=2^m+\sum_{j=1}^m s_j2^{m-j}.
\end{equation}
Strings of length $m$ use the integers from $2^m$ through $2^{m+1}-1$,
so this is an injective coding of all nonempty finite binary strings.
For an infinite binary sequence $\alpha=(\alpha_1,\alpha_2,\ldots)$, set
\begin{equation}\label{tail:eq:Aalpha}
 A_\alpha=\{N(\alpha_1,\ldots,\alpha_m):m\geq1\}.
\end{equation}
Each $A_\alpha$ is infinite. Two different sequences share only the
codes of their common initial prefixes, hence have finite intersection.
In a finite collection of distinct sequences, each corresponding set
therefore has infinitely many indices belonging to no other set in the
collection. We call these its private indices.

\begin{theorem}[Explicit maximal differential transcendence]\label{tail:thm:surreal}
Put $\xi_\alpha=\xi_{A_\alpha}$ using~\eqref{tail:eq:xiA} and~\eqref{tail:eq:Aalpha}.
Then
\begin{equation}\label{tail:eq:allrealjets}
 \{\BM^k\xi_\alpha:
       \alpha\in\{0,1\}^{\N},\ k\geq0\}
\end{equation}
is algebraically independent over $B_0=\Q((t^\Q))$.
In particular, with the induced Berarducci--Mantova derivation,
\begin{equation}\label{tail:eq:maxdtrdeg}
 \dtrdeg_{B_0}H_0=\trdeg_{B_0}H_0=\cf.
\end{equation}
The map sending free ordinary differential indeterminates $X_\alpha$
to $\xi_\alpha$ embeds the differential polynomial algebra
$B_0\{X_\alpha:\alpha\in\{0,1\}^{\N}\}$ into $H_0$.
\end{theorem}
\begin{proof}
It suffices first to prove independence of the Euler jets belonging to
finitely many distinct sequences $\alpha_1,\ldots,\alpha_r$, up to order
$K$. Their combined coefficient vector at $n$ is
\[
 v_n=\bigl(\mathbf1_{A_{\alpha_a}}(n)n^k
                  \bigr)_{1\leq a\leq r,\ 0\leq k\leq K}
       \in\Q^{r(K+1)}.
\]
On the infinitely many private indices of $A_{\alpha_a}$, all blocks
except block $a$ are zero. Within that block the vectors are Vandermonde
vectors. Every cofinite subfamily of these private indices spans the
whole $a$th block. Doing this for all $a$ shows that the cofinite span
of the full combined family is $\Q^{r(K+1)}$. Theorem~\ref{tail:thm:tail}
therefore proves independence of the combined finite jet tuple.
The triangular change~\eqref{tail:eq:BMtriangular} gives the same result for
$\BM$-jets. Since any polynomial involves only finitely many sequences
and jet orders,~\eqref{tail:eq:allrealjets} follows.

There are $\cf$ binary sequences and the $\xi_\alpha$ are pairwise
distinct, so this gives differential transcendence degree at least $\cf$.
The matching upper bound is a separate cardinality count that uses none of
the independence just proved: $\AR$ and $\Q$ are countable, and a Hahn
series in $H_0$ is in particular a function from $\Q$ to the countable set
$\AR$, so $|H_0|\leq\aleph_0^{\aleph_0}=\cf$. A transcendence degree is
bounded by the cardinality of the field, so both displayed degrees are at
most $\cf$; the exhibited family supplies the reverse inequality for the
differential one, and a differentially independent family is in particular
algebraically independent. This also gives $|H_0|=\cf$. Finally, a differential polynomial
is an ordinary polynomial in finitely many formal jets. Its evaluation
is injective precisely by the proved jet independence.
\end{proof}

\begin{remark}[The index family is not arbitrary]
\label{tail:rem:indexfamily}
Theorem~\ref{tail:thm:surreal} is about the binary-prefix family
$(A_\alpha)$ of~\eqref{tail:eq:Aalpha}, not about all subsets of
$\N_{>0}$. The unrestricted family $(\xi_A)_{A\subseteq\N_{>0}}$ is
\emph{not} independent: for any $A$ and $B$,
$\mathbf1_{A\cup B}+\mathbf1_{A\cap B}=\mathbf1_A+\mathbf1_B$, so
\[
 \xi_{A\cup B}+\xi_{A\cap B}=\xi_A+\xi_B ,
\]
already a $\Q$-linear relation. Almost disjointness is what removes such
relations, by giving each member infinitely many private indices; the same
role is played in the proof by the cofinite span. Theorem~\ref{tail:thm:onesurreal}
is the complementary statement that no hypothesis at all beyond infinitude
is needed when only \emph{one} $A$ is considered.
\end{remark}

Two reports placed later bear on this remark from other directions.
\path{docs/surreal/transcendence-over-bounded-support/} proves independence
over the fraction field of bounded-support series, from the support rather
than the coefficients (\texttt{bst:sec:collection}); its integer-coefficient
witnesses lie in $B_0$, and its family indexed by \emph{all} subsets is
independent because its coefficients are finite-pattern codes.
The report on hidden negative Hermitian directions uses
prime-denominator tails $\sum_pt^{1-1/p}$ as the exponent-side counterpart
over its finite-lattice-supported base (\texttt{hnd:thm:independence},
\texttt{hnd:cor:degreebound}); for $\Gamma=\Q$ the two bases do not contain
each other.

\emph{Related (batch 34).} The coefficient-field part of
\path{docs/surreal/independent-surreal-copies/} reuses the binary-prefix
family~\eqref{tail:eq:Aalpha}. Its \texttt{isc:cp:cor:continuum} gives $\cf$
elements $\sum_{n\in A_\alpha}x^{n!}t^n$ of $k(x)((t))$, with $x$
transcendental and $k$ any field of characteristic zero, that are
differentially algebraically independent over $k((t))(x)$ for the Euler
derivation $t\,d/dt$. The mechanism is \texttt{isc:cp:thm:lacunary}, the
rapid growth of the degree in $x$ of late coefficients, not a tail span.
\texttt{isc:cp:thm:generic} and \texttt{isc:cp:cor:tree} put algebraically
independent coefficients at private indices instead, and use coefficient
derivations and a Vandermonde argument, as in the proof of
Theorem~\ref{tail:thm:surreal}. When the coefficient field has transcendence
degree at least $\kappa$, they give $\kappa^{\aleph_0}$ series independent
over a finite-generation coefficient envelope. Those bases and derivations
are not $B_0$ and $\BM$, so neither result is deduced from
Theorem~\ref{tail:thm:surreal} or implies it.

\emph{Related (batch 35).} The Hahn-join part of the same report (labels
\texttt{isc:hj:}) uses the branch sets~\eqref{tail:eq:Aalpha} shifted down
by one, that is, the code $N(s)-1$ of~\eqref{tail:eq:binarycode} on
$0$-based indices (\texttt{isc:hj:eq:branchsets},
\texttt{isc:hj:lem:branches}). With primes numbered from $p_0=2$ and
$\rho_n=p_{2n}^{-3/2}+p_{2n+1}^{-3/2}$, its \texttt{isc:hj:main:prime}
indexes by them $\cf$ Hahn sums $1+\sum_n\omega^{\rho_n}$, each over one
shifted branch set, that are algebraically independent over the compositum
of two Archimedean full Hahn fields, those with exponent groups spanned by
the $p_{2n}^{-3/2}$ and by the $p_{2n+1}^{-3/2}$. The linear independence of
square roots of distinct primes enters there in the exponents, since
$p^{-3/2}=\sqrt p/p^2$, rather than in the coefficients as
in~\eqref{tail:eq:xiA}. The mechanism is a mixed-support theorem
(\texttt{isc:hj:main:mixed}), proved with derivations from coordinate
characters of the two exponent groups and a finite differential annihilator
(\texttt{isc:hj:lem:annihilator}), not a tail span; its independence is
algebraic rather than differential, and its base is not $B_0$.

\subsection{Why the base field and the derivation matter}
The theorem is not merely independence over $\Q(\omega)$. The base
$B_0$ contains \emph{every} rational-coefficient Hahn series with rational
valuation support, including supports with infinitely many exponents
below a finite rational bound. Nevertheless the displayed numbers escape
all finite algebraic differential equations over that entire field.
Conversely, independence over $B_0$ does not mean independence over the
whole $H_0$, which already contains them.

Let $\mathcal A_0$ be the relative algebraic closure of $B_0$ in $H_0$.
By Theorem~\ref{tail:thm:coefficient}, it is the union of
$F((t^\Q))$ over real number fields $F\subseteq\AR$. It is stable under
$\BM$. The same jet family is independent even over $\mathcal A_0$.
Indeed, adjoining an algebraic base extension does not change the
transcendence degree of a finite tuple: the compositum is algebraic over
the field of that tuple. Thus
\[
 \dtrdeg_{\mathcal A_0}H_0=\cf.
\]
In particular, every individual coefficient $\sqrt{p_n}$ is already in
the algebraic base, but the infinite sum is not differentially algebraic
there. Algebraic closure and unrestricted Hahn formation are decisively
different operations.

\section{Analytic coefficient fields and two commuting derivations}\label{tail:sec:analyticsetup}

\subsection{A common disk and independent quadratic functions}
Let $M=\C(z)$ and, for $n\geq1$, define
\begin{equation}\label{tail:eq:hn}
 a_n(z)=1-\frac z{2^n},\qquad
 h_n(z)=\sqrt{1-z/2^n},\qquad h_n(0)=1.
\end{equation}
Each branch is holomorphic on the disk $D_2=\{|z|<2\}$; the first
possible branch point is at $2$. We usually work on the smaller common
disk $D=\{|z|<1\}$, which also accommodates the dilation identity below.
The field
\[
 L=\C(z)(h_n:n\geq1)
\]
is realized inside the field $\mathscr M_0$ of meromorphic germs at $0$.
The common ambient Hahn field $\mathscr M_0((t^\Q))$ contains both
$L((t^\Q))$ and the ring $\OO(D)((t^\Q))$; neither latter object is
asserted to contain the other. This choice of germs fixes compatible square-root choices; finite rational
combinations are meromorphic germs even when they do not have a common
global domain.

\begin{lemma}[Independent analytic square classes]\label{tail:lem:analyticclasses}
The $a_n$ have independent square classes in $\C(z)^\times/\C(z)^{\times2}$.
\end{lemma}
\begin{proof}
For any nonempty finite $J$, the product $\prod_{n\in J}(1-z/2^n)$ has
a simple zero at each $2^n$ with $n\in J$. A square rational function has
even order at every point, since zeros and poles of its square have twice
the original integer orders. The product is therefore not a square.
\end{proof}

Our ambient and base fields are
\begin{equation}\label{tail:eq:analyticfields}
 H=L((t^\Q)),\qquad B=\C(z)((t^\Q)),\qquad K_\Q=\C((t^\Q)).
\end{equation}
The smaller rational function field $K_\Q(z)$ embeds in $B$. Indeed,
the map from $K_\Q[z]$ is injective by comparing Hahn coefficients:
a polynomial relation in $z$ would make every ordinary complex polynomial
coefficient vanish. It therefore extends to the fraction field.

\begin{warning}[A large relation field is not a common-domain ring]
An arbitrary element of $B=\C(z)((t^\Q))$ can have coefficient poles
approaching $0$. Hence $B$ is not being identified with a ring of
functions on one common disk. Algebraic relations over $B$ are interpreted
in the formal Hahn field $H$ with meromorphic-germ coefficients. The
specific functions constructed here \emph{do} lie in
$\OO(D)((t^\Q))$. A passage to actual common-domain or halo functions is
made only where justified in Section~\ref{tail:sec:halos}. This distinction
agrees with the ring distinctions in the repository catalogue
\cite{RepoMap}.
\end{warning}

\subsection{The derivations and their exact coefficient formula}
The ordinary derivation $\partial=\partial_z$ of $\C(z)$ extends to $L$:
by implicit differentiation of $h_n^2=a_n$,
\[
 \partial h_n=\frac{h_n}{2(z-2^n)}\in L.
\]
Define on $H$
\begin{equation}\label{tail:eq:twoderivations}
 \partial\left(\sum_q f_qt^q\right)=\sum_q f_q't^q,
 \qquad
 \EDer\left(\sum_q f_qt^q\right)=\sum_q qf_qt^q.
\end{equation}
Both maps preserve supports up to deletion. Coefficientwise finite
convolution proves the Leibniz rule, and coefficient comparison proves
that they commute. They preserve $B$ as well. The same formulas also define commuting derivations on
$\mathscr M_0((t^\Q))$ and preserve its common-domain holomorphic
coefficient subring $\OO(D)((t^\Q))$.

Let
\[
 c_j=\ff{1/2}{j}=\frac12\left(\frac12-1\right)\cdots
                         \left(\frac12-j+1\right),\qquad c_0=1.
\]
Every $c_j$ is nonzero. Repeated differentiation gives
\begin{equation}\label{tail:eq:hnjet}
 \partial^jh_n(z)=c_j(z-2^n)^{-j}h_n(z).
\end{equation}
For example, differentiating the right side at order $j$ multiplies it
by $(1/2-j)/(z-2^n)$, giving the next falling factorial.

For every subset $A\subseteq\N_{>0}$, set
\begin{equation}\label{tail:eq:FA}
 F_A(z,t)=\sum_{n\in A}t^nh_n(z)\in H\cap\OO(D)((t^\Q)).
\end{equation}
The support is a subset of the positive integers, so this is a Hahn
series without any coefficient growth condition. Its mixed jets have the
particularly useful form
\begin{equation}\label{tail:eq:mixedcoefficients}
 \partial^j\EDer^kF_A
   =\sum_{n\in A}t^nh_n(z)\,
             c_j\frac{n^k}{(z-2^n)^j}.
\end{equation}

\section{Every mixed jet is algebraically independent}\label{tail:sec:mixed}

\subsection{An elementary exponential-graph lemma}
\begin{lemma}\label{tail:lem:expgraph}
If $R(X,Y)\in\C[X,Y]$ is nonzero, then $R(n,2^n)=0$ for only finitely
many positive integers $n$.
\end{lemma}
\begin{proof}
Write $R(X,Y)=\sum_{m=0}^M p_m(X)Y^m$, with $p_M\ne0$.
If $M=0$, this is the fact that a nonzero polynomial has finitely many
roots. In general, for sufficiently large real $n$, $p_M(n)\ne0$ and
\[
 \frac{R(n,2^n)}{p_M(n)2^{Mn}}
   =1+\sum_{m<M}\frac{p_m(n)}{p_M(n)}\,2^{-(M-m)n}.
\]
Each rational function of $n$ in the sum has at most polynomial growth,
while its exponential factor tends to zero. The displayed ratio tends
to $1$, so it is nonzero for all sufficiently large $n$. Only finitely
many positive integers remain.
\end{proof}

The lemma is an asymptotic statement about an ordinary finite polynomial;
it invokes no transcendence conjecture and no assertion about the
arithmetic nature of individual values $2^n$.

\subsection{Every infinite subfamily spans every finite jet space}
\begin{lemma}[Hereditary full jet span]\label{tail:lem:jetspan}
Fix $J,K\geq0$. For $n\geq1$, form the vector
\[
 v_n(z)=\left(c_j\frac{n^k}{(z-2^n)^j}
          \right)_{0\leq j\leq J,\ 0\leq k\leq K}
       \in M^{(J+1)(K+1)},\qquad M=\C(z).
\]
Every infinite subfamily of these vectors spans the entire displayed
space over $M$. In particular, every cofinite subfamily of every infinite
subfamily has full span.
\end{lemma}
\begin{proof}
Suppose an infinite index set $A$ failed to span. A nonzero $M$-linear
functional would annihilate all its vectors. Absorb the nonzero $c_j$
into its coefficients and clear their common rational denominator. There
would exist polynomials $C_{jk}(z)\in\C[z]$, not all zero, such that
\begin{equation}\label{tail:eq:jetlinearrelation}
 \sum_{j,k} C_{jk}(z)\frac{n^k}{(z-2^n)^j}=0
 \quad\text{in }\C(z),\quad n\in A.
\end{equation}
Choose the largest $j_*$ for which some $C_{j_*k}$ is nonzero. Multiply
by $(z-2^n)^{j_*}$ and evaluate the resulting polynomial identity at
$z=2^n$. This evaluates a cleared polynomial identity, not the analytic series at
a branch point. All lower-$j$ terms vanish, giving
\[
 \sum_{k=0}^K C_{j_*k}(2^n)n^k=0\qquad(n\in A).
\]
The polynomial
$R(X,Y)=\sum_{k=0}^K C_{j_*k}(Y)X^k$ is nonzero. This contradicts
Lemma~\ref{tail:lem:expgraph}, because $A$ is infinite. The same proof covers
$j_*=0$: the polynomial identity itself may simply be evaluated.
\end{proof}

\subsection{The mixed-jet theorem}
\begin{theorem}[Hereditary mixed differential independence]\label{tail:thm:mixed}
For every infinite $A\subseteq\N_{>0}$, the family
\begin{equation}\label{tail:eq:mixedindependence}
 \{\partial^j\EDer^kF_A:j,k\geq0\}
\end{equation}
is algebraically independent over $B=\C(z)((t^\Q))$.
Equivalently, the homomorphism from the free differential polynomial
algebra in one indeterminate with two commuting derivations to $H$,
sending that indeterminate to $F_A$, is injective.
\end{theorem}
\begin{proof}
Any finite collection of mixed jets is contained in a rectangle
$0\leq j\leq J$, $0\leq k\leq K$. By~\eqref{tail:eq:mixedcoefficients}, the
vector of all jets in the rectangle has the form of
Theorem~\ref{tail:thm:tail}, with $r_n=h_n$, $\gamma_n=n$, and coefficient
vectors from Lemma~\ref{tail:lem:jetspan}. Lemma~\ref{tail:lem:analyticclasses}
supplies square-class independence, and Lemma~\ref{tail:lem:jetspan} supplies
full cofinite span. Thus $E=\varnothing$ and the entire finite rectangle
is algebraically independent over $B$. This proves the infinite family
assertion. The formal mixed jets in a differential polynomial algebra
with commuting derivations are indexed by pairs $(j,k)$; their evaluation
is injective exactly when the family~\eqref{tail:eq:mixedindependence} is
algebraically independent.
\end{proof}

This excludes every finite-order algebraic partial differential equation
with coefficients in $B$, not just linear equations, autonomous equations,
or equations of some bounded order. Setting $k=0$ proves ordinary
differential transcendence in $z$; setting $j=0$ proves differential
transcendence for the scale derivation. Because $K_\Q(z)\subseteq B$,
the result is stronger than independence over the rational function field
with surcomplex coefficients from $K_\Q$.

\begin{remark}[The upgrade a companion warning declines]
\label{tail:rem:dynupgrade}
The companion report on dynamics and normal forms~\cite{RepoDyn} proves
an algebraicity dichotomy for its Schr\"oder-type solutions
$H_\varepsilon$: for $V\in\C[w]$ with $V(0)=0$ and $V'(0)\neq0$, the
solution $H_\varepsilon$ is algebraic over
$\mathscr B=\C(w)((t^\Gamma))$ if and only if $\deg V=1$, so it is
transcendental over $\mathscr B$ whenever $\deg V\geq2$.
Its warning \texttt{dyn:warn:diffalg}
deliberately refuses to upgrade that to ``differentially
transcendental'', because its own $H_\varepsilon$ satisfies the first-order
linear equation $H_\varepsilon'=\Omega_\varepsilon H_\varepsilon$ with $\Omega_\varepsilon\in\mathscr B$.
That refusal is correct and is not weakened by anything here: a function
with a first-order linear equation over the base is differentially
algebraic over it, whatever else is true.

Theorem~\ref{tail:thm:mixed} supplies the upgrade for \emph{different}
functions over a base of the same shape. With $k=0$ it says that
$\{\partial_z^jF_A:j\geq0\}$ is algebraically independent over
$\C(z)((t^\Q))$, that is, $F_A$ satisfies no algebraic differential
equation of any finite order over that field --- the exact conclusion
\texttt{dyn:warn:diffalg} declines for $H_\varepsilon$. The base is the
$\Gamma=\Q$ case of the same $\C(z)((t^\Gamma))$ shape, and by
Theorem~\ref{tail:thm:continuumanalytic} the statement holds for
continuum many functions on the one fixed disk $D=\{|z|<1\}$, with no
shrinking of $D$ as the parameter varies. The two reports therefore
describe different functions, not a disagreement: the companion's
$H_\varepsilon$ remains differentially algebraic over $\mathscr B$, and $F_A$ is not.
The contrast is sharpened rather than softened by
\eqref{tail:eq:dilation} below, where an $F_A$ that satisfies a clean
dilation functional equation is nevertheless differentially
transcendental.
\end{remark}

As in the actual-surreal application, the same finite jet families remain
independent after any algebraic base extension of $B$ in the ambient
field. This follows from invariance of transcendence degree under
algebraic base extension, not from a claim that the original sign
changes fix that enlarged base.

\subsection{A functional equation does not supply a differential equation}
Take $F=F_{\N_{>0}}$. Reindexing a Hahn-summable family gives the exact
identity
\begin{equation}\label{tail:eq:dilation}
 F(2z,t)-tF(z,t)=t\sqrt{1-z},\qquad |z|<1.
\end{equation}
Indeed, the $n=1$ term of $F(2z,t)$ is $t\sqrt{1-z}$, and its remaining
terms are $tF(z,t)$. The expression $F(2z,t)$ is legitimate here because
the coefficients of $F$ are holomorphic on $D_2$.

Thus a simple dilation relation and quadratic coefficient functions
coexist with complete mixed-jet algebraic independence. Dilation is an
additional operator; replacing it by a finite collection of derivatives
would be an invalid inference. The result does not exclude functional
equations involving shifted or dilated arguments.

\subsection{Rational Taylor data and a summation warning}
From~\eqref{tail:eq:hnjet}, coefficientwise Taylor extraction yields
\begin{equation}\label{tail:eq:taylorrational}
 \partial^jF(0,t)=(-1)^jc_j\sum_{n\geq1}2^{-nj}t^n
          =(-1)^jc_j\frac{t}{2^j-t}\in\C(t).
\end{equation}
Every specialized mixed jet $\partial^j\EDer^kF(0,t)$ is therefore
rational in $t$ as well. In the ordinary formal power-series ring
$K_\Q[[z]]$, the Taylor image of $F$ is
\begin{equation}\label{tail:eq:formaltaylor}
 \sum_{j\geq0}(-1)^j\binom{1/2}{j}
                       \frac{t}{2^j-t}\,z^j.
\end{equation}
To verify this equality, fix $j$, extract the ordinary complex Taylor
coefficient of each $h_n$, and then sum the Hahn family indexed by $n$.
This is a coefficient-by-coefficient construction of a formal $z$-series.

\begin{warning}\label{tail:warn:taylor}
For an ordinary nonzero complex value of $z$, expression
\eqref{tail:eq:formaltaylor} is \emph{not} a Hahn sum indexed by $j$:
infinitely many of its terms contribute to exponent $t^1$. Its meaning
is as a formal Taylor series in $z$, or as a jointly summable expression
when $z$ is a suitable infinitesimal. One must not interchange the two
infinite summation notions without a support certificate.
\end{warning}

There is also no contradiction between~\eqref{tail:eq:taylorrational} and
Theorem~\ref{tail:thm:mixed}. The theorem is about functions in $z$ over the
specified function field. Evaluation at a point need not preserve
algebraic independence. In fact, evaluation at any ordinary $z_0\in D$
places these values inside $K_\Q$ itself. Unlike the actual numbers of
Section~\ref{tail:sec:surreal}, their point values are not being asserted
transcendental over that field.

\section{A continuum of jointly free surcomplex analytic jets}\label{tail:sec:analyticcontinuum}

Use the almost disjoint sets $A_\alpha$ from~\eqref{tail:eq:Aalpha} and define
\[
 G_\alpha(z,t)=F_{A_\alpha}(z,t)
       =\sum_{n\in A_\alpha}t^n\sqrt{1-z/2^n}.
\]
Every $G_\alpha$ has the same common holomorphic coefficient disk $D$.
No shrinking of the domain is required as $\alpha$ varies.

\begin{theorem}[Continuum many jointly independent mixed jets]
\label{tail:thm:continuumanalytic}
The family
\begin{equation}\label{tail:eq:allanalyticjets}
 \{\partial^j\EDer^kG_\alpha:
       \alpha\in\{0,1\}^{\N},\ j,k\geq0\}
\end{equation}
is algebraically independent over $B=\C(z)((t^\Q))$.
It embeds the free differential polynomial algebra on $\cf$ generators,
with two commuting derivations, into $H=L((t^\Q))$.
Moreover $\trdeg_B H=\cf$; the constructed family has the largest
possible cardinality for algebraic independence in this workspace.
\end{theorem}
\begin{proof}
Choose finitely many distinct $\alpha_1,\ldots,\alpha_r$ and a finite
rectangle of jet orders. At each coefficient index $n$, concatenate the
jet coefficient vector of Lemma~\ref{tail:lem:jetspan} in block $a$ if
$n\in A_{\alpha_a}$, and use a zero vector in that block otherwise.
For each $a$, its infinite private index set contributes only to block
$a$. After any finite deletion, Lemma~\ref{tail:lem:jetspan} shows that these
vectors span the entire block. The cofinite span of the combined vectors
is consequently the full combined jet space. Theorem~\ref{tail:thm:tail}
proves algebraic independence of the chosen finite tuple, and hence of
all of~\eqref{tail:eq:allanalyticjets}.

A differential polynomial uses only finitely many generators and mixed
jets, giving the claimed differential-algebra embedding. Finally,
$|\C(z)|=\cf$, the algebraic extension $L$ has cardinality $\cf$, and
$\Q$ is countable. Thus $|H|\leq\cf^{\aleph_0}=\cf$. The exhibited
independent family has cardinality $\cf$, proving the exact
transcendence degree and the maximality assertion.
\end{proof}

\begin{remark}[No assertion of a transcendence basis]
The cardinal equality does not say that the displayed jet family is a
transcendence \emph{basis}, or that $H$ is algebraic over the field it
generates. It says that the family is independent and has maximal possible
cardinality. Proper subsets of an infinite basis can have the same
cardinality as the basis, so cardinal maximality must not be confused
with maximality under inclusion.
\end{remark}

\Needspace{7\baselineskip}
\section{From coefficient functions to surcomplex halos}\label{tail:sec:halos}

\subsection{Support-certified infinitesimal substitution}
An actual surcomplex halo value may involve scales beyond $\Q$. Let
$\Gamma$ be a set-sized ordered subgroup of the surreal exponent group
containing $\Q$, and put $K_\Gamma=\C((t^\Gamma))\subseteq\No[i]$ via
$t^\gamma=\omega^{-\gamma}$. Let $z_0\in D$ be an ordinary complex
number and let $\varepsilon\in K_\Gamma$ be infinitesimal, meaning either
$\varepsilon=0$ or $\supp(\varepsilon)\subset\Gamma_{>0}$.
Define
\begin{equation}\label{tail:eq:halovalue}
 F_A(z_0+\varepsilon,t)
  =\sum_{n\in A}\sum_{\ell\geq0}
       t^n\frac{h_n^{(\ell)}(z_0)}{\ell!}\varepsilon^\ell.
\end{equation}

\begin{proposition}[Joint summability of the halo formula]\label{tail:prop:halo}
Formula~\eqref{tail:eq:halovalue} is a Hahn-summable double family and gives
an element of $K_\Gamma$. The same holds for every mixed jet and every
finite polynomial combination of the functions considered above.
\end{proposition}
\begin{proof}
The case $\varepsilon=0$ reduces to the original sum. Otherwise write
$T=\supp(\varepsilon)\subset\Gamma_{>0}$ and $S=A\subset\N_{>0}$.
By the positive-support Neumann lemma, the monoid $T^*$ generated by
$T$, including zero, is well ordered. For each exponent it has only
finitely many representations as a word in $T$, counting all possible
word lengths. Therefore the family $(\varepsilon^\ell)_{\ell\geq0}$
is Hahn summable. The set $S+T^*$ is well ordered, with finitely many
pairs contributing at any fixed exponent. These two finiteness
statements together give joint summability of the $(n,\ell)$ family.
The complex Taylor coefficients do not enlarge any support. Derivatives
of $h_n$ are holomorphic at the same $z_0$, and finite products use the
same support lemmas, proving the remaining assertions.
\end{proof}

This is a direct certificate in an explicitly chosen set-sized Hahn
workspace. It is not a convergence assertion for partial sums in the
fine topology of the entire class $\No[i]$, and it makes no claim that
$F_A$ is defined at every infinite surcomplex argument.

\subsection{Functional independence is visible on ordinary points}
Although the largest coefficient field $B$ is a formal meromorphic-germ
Hahn field, the independence conclusion has an ordinary-point functional
meaning over $K_\Q(z)$.

\begin{proposition}\label{tail:prop:functionalmeaning}
No nonzero polynomial over $K_\Q(z)$ in finitely many of the mixed jets
of Theorem~\ref{tail:thm:continuumanalytic} vanishes identically on their
common ordinary disk of definition. The same polynomial cannot vanish
identically on all the corresponding halo values.
\end{proposition}
\begin{proof}
Clear the finitely many rational-function denominators, obtaining a
nonzero polynomial with coefficients in $K_\Q[z]$. Its evaluation on the
chosen jets is a nonzero element of $H$ by
Theorem~\ref{tail:thm:continuumanalytic}. In fact it lies in
$\OO(D)((t^\Q))$: only finite products and sums of common-domain
coefficient functions have been used. At least one of its Hahn
coefficients is therefore a nonzero holomorphic function on $D$.
It is nonzero at some ordinary point, and indeed outside its discrete
zero set.

For each denominator in $K_\Q[z]$, the coefficient of its lowest
valuation is a nonzero ordinary complex polynomial. Away from the finite
union of their zeros, the denominator does not specialize to zero.
Choose a point where the numerator coefficient is nonzero and these
finitely many denominator coefficients are nonzero. The original
polynomial expression is nonzero there. Ordinary points are halo points
with $\varepsilon=0$, proving the last assertion.
\end{proof}

\subsection{Three differentiations that must not be conflated}
In Section~\ref{tail:sec:surreal}, $\BM$ acts on actual surreal numbers and
satisfies $\BM=-t\EDer$ on the stated rational-monomial fields. In the
analytic coefficient field, $\partial_z$ differentiates the independent
complex parameter, while $\EDer$ acts on the Hahn scale and fixes the
coefficient functions. These are two explicitly defined commuting
operators on a function field.

Specializing $z$ to a nonconstant surreal expression can introduce a
chain-rule term under $\BM$. It is therefore incorrect to identify
$\partial_z$ with $\BM$, or to assert without qualification that the
function-field derivations remain those same operators after every
surreal substitution. The algebraic independence theorem is formulated
before such substitutions; Section~\ref{tail:sec:surreal} separately supplies
a direct, normalized Berarducci--Mantova result for actual numbers.

\section{Algebraic approximation depends on the value group}\label{tail:sec:approx}

Theorem~\ref{tail:thm:coefficient} also gives a sharp approximation criterion.
It explains how a transcendental series can be approximated arbitrarily
closely by algebraic ones in one workspace, but not after the value group
is enlarged.

\begin{theorem}[Closure of the algebraic coefficient part]\label{tail:thm:closure}
Let $B=M((t^\Gamma))$, let $\overline M$ be an algebraic closure in
characteristic zero, and let $\mathcal A$ be the relative algebraic
closure~\eqref{tail:eq:relativeunion}. A series
$f=\sum_\eta c_\eta t^\eta\in\overline M((t^\Gamma))$ lies in the
valuation-topological closure of $\mathcal A$ if and only if
\begin{equation}\label{tail:eq:prefixcriterion}
 [M(c_\eta:\eta\leq\gamma):M]<\infty
                    \qquad\text{for every }\gamma\in\Gamma.
\end{equation}
Zero coefficients may be omitted from the generated field.
\end{theorem}
\begin{proof}
The valuation neighborhoods of $f$ are the sets
$\{g:v(f-g)>\gamma\}$ for $\gamma\in\Gamma$.
If $g\in\mathcal A$ belongs to this neighborhood, its coefficients
through exponent $\gamma$ agree with those of $f$. All coefficients of
$g$ lie in one finite extension of $M$ by
Theorem~\ref{tail:thm:coefficient}, proving necessity.

Conversely, suppose~\eqref{tail:eq:prefixcriterion}. The truncation
$f_{\leq\gamma}=\sum_{\eta\leq\gamma}c_\eta t^\eta$ has well-ordered
support and coefficients in the finite extension from that condition.
It lies in $\mathcal A$, and either equals $f$ or its difference from $f$
has least exponent strictly greater than $\gamma$. It approximates $f$
in the required neighborhood, proving sufficiency.
\end{proof}

For $\Gamma=\Q$, the series $\sum_{n\geq1}\sqrt{p_n}t^n$ satisfies the
criterion: below any rational cutoff there are only finitely many of its
positive integer exponents. It is therefore a valuation limit of its
finite algebraic truncations, despite its differential transcendence.

For an enlargement, take $\Gamma=\Q^2$ in lexicographic order,
embedding the old value group by $q\mapsto(0,q)$, and put
\[
 f=\sum_{n\geq1}\sqrt{p_n}\,t^{(0,n)}.
\]
Its support is still well ordered and has order type $\omega$. However,
every one of its exponents lies below $(1,0)$. At that cutoff the
coefficient field is the infinite multiquadratic extension, so the
criterion fails. No element algebraic over the enlarged base
$M((t^{\Q^2}))$, here with $M=\Q$, approximates $f$ with error valuation
greater than $(1,0)$. In particular this is an obstruction to approximation
by \emph{any} algebraic element, not just by finite truncations.

Under the embedding $(a,b)\mapsto a\omega+b$ of value groups, the
second series has the same displayed surreal expansion
$\sum_n\sqrt{p_n}\omega^{-n}$. The enlarged workspace has an additional
precision scale $\omega^{-\omega}$ which no finite truncation reaches.
Thus the topology and the available precision scales must be specified;
Hahn summability alone is not valuation convergence of finite truncations
in an arbitrary-rank field.

\section{Effective limitations and exact verification}\label{tail:sec:effective}

\subsection{Algebraicity is undecidable for unrestricted computable streams}
Consider the representation
\[
 x_\beta=\sum_{n\geq1}\beta_n\sqrt{p_n}\,t^n,
 \qquad \beta_n\in\{0,1\},
\]
where a program is promised to compute the total sequence $(\beta_n)$.
Every coefficient has a finite exact description; the support is
contained in the positive integers. Hence no support-validity problem
is hidden in this representation.

\begin{proposition}\label{tail:prop:undecidable}
There is no algorithm that, for every such total computable coefficient
stream, decides whether $x_\beta$ is algebraic over $B_0$.
There is also no algorithm deciding whether it is differentially
algebraic over $B_0$ for $\BM$.
\end{proposition}
\begin{proof}
Given a program $e$ on a fixed input, define $\beta_n$ by simulating $e$
for $n$ steps: return $1$ if it has not yet halted and $0$ otherwise.
This coefficient program is total regardless of whether $e$ halts.
If $e$ halts, only finitely many $\beta_n$ are nonzero, so $x_\beta$ is
algebraic and therefore differentially algebraic. If $e$ never halts,
$x_\beta=\xi_{\N_{>0}}$, which is differentially transcendental by
Theorem~\ref{tail:thm:onesurreal}, and in particular transcendental.
Either proposed decision procedure would decide whether arbitrary
programs halt. Such a procedure is impossible by the standard diagonal
argument: a program that halts exactly when the purported halting test
says it does not halt contradicts the test applied to its own code.
\end{proof}

Every finite set of coefficient observations is also compatible with
both a finite-support algebraic continuation and an infinite-support
differentially transcendental continuation. The latter assertion uses
the fact that \emph{every} infinite subset works, not just the full set
of positive integers. Thus the finite certificates of
Corollary~\ref{tail:cor:finitewitness} require structural information about the
stream; numerical or symbolic inspection of a long prefix is not a
substitute for such information.

\subsection{What the accompanying checks do}
The supplied \path{code/verify.py} uses exact arithmetic and symbolic
identities. It checks the top-degree finite-difference formula, small
multiquadratic sign products, the quadratic exceptional relation,
Vandermonde and mixed-jet ranks, the derivative formulas, the finite
boundary term in the dilation identity, Taylor coefficient identities,
and the finite behavior of the binary-prefix construction.
All 404 recorded checks passed with Python 3.13.5 and SymPy 1.14.0;
the exact environment is also recorded in
\path{data/verification_results.json}.

These checks are deliberately subordinate to the proofs. A finite rank
calculation does not prove the hereditary full-span lemma for all orders
and all infinite subsets; that lemma is proved by the pole argument and
Lemma~\ref{tail:lem:expgraph}. Finite sign-orbit examples do not prove the
classification for arbitrary well-ordered supports; that theorem is
proved by exact Galois actions, finite differences, and cofinite spanning.
The existing module \path{Surreal/Algebra/TailSpan.lean} proves
Lemmas~\ref{tail:lem:stabilization} and~\ref{tail:lem:multilinear}.
The precise scope is recorded in \path{docs/FORMALIZATION.md}; the other
theorems in this report are not thereby formalized.

\subsection{Proof dependencies and falsifiable checkpoints}
The core proof has a short dependency structure. Finite multiquadratic
Galois theory gives independently available sign changes. A degree-$D$
polynomial relation then forces its highest polarization to vanish on
every distinct-index $D$-tuple. Cofinite spanning forces that polarization
to vanish identically. Finally, one transverse linear functional recovers
exactly the finite exceptional multiquadratic extension. Each of these
steps is isolated in a separately stated lemma or theorem.

The two applications add different spanning arguments. Actual surreal
Euler jets use the Vandermonde determinant. Analytic mixed jets use the
largest pole order and the impossibility of a nonzero polynomial vanishing
on infinitely many points $(n,2^n)$. The continuum constructions add only
the finite-overlap property of binary-prefix sets.

The assumptions have specific roles. Distinct support exponents ensure
that sign-conjugate primitive elements are distinguishable coefficientwise.
Independent square classes license independent sign changes. Characteristic
zero ensures that the top polarization is detected by $D!$ and that
separability causes no coefficient-orbit pathology. Finite-dimensional
vector spaces make the cofinite span stabilize after finitely many
deletions. Dropping any of these hypotheses requires a new argument,
not a notational adjustment.

\section{Further questions and final assessment}\label{tail:sec:questions}

The article establishes an exact relation theory for one substantial
class of Hahn constructions. It does not classify arbitrary algebraic
coefficient families. When the available coefficient Galois actions are
correlated rather than independent sign changes, their orbits need not
supply the independent translation directions used in the proof.
A natural next problem is to replace the square-root hypotheses with a
usable condition on finite Galois representations, and to determine which
parts of the affine-plane classification survive.

A second problem is effectivity under certificates. For a stream supplied
with a finite-state recurrence, a certified pole pattern, or an explicit
cofinite-spanning proof, one can ask for algorithms that compute its
transcendence dimension and exceptional field. The unrestricted decision
problem of Proposition~\ref{tail:prop:undecidable} does not rule out such
structured subclasses. Corollary~\ref{tail:cor:finitewitness} suggests a
bounded-degree certificate format based on disjoint spanning sets and a
nonzero mixed difference.

A third problem concerns the interaction with richer surreal derivations
and larger monomial groups. The actual-number theorem here deliberately
uses rational exponents, where the normalized Berarducci--Mantova
restriction is explicit. Extending its transparent jet-spanning proof to
supports involving general surreal exponents would require control of
their derivatives, not just their order. Likewise, the analytic theorem
does not automatically survive substitutions in which the parameter
itself carries a nontrivial surreal derivative.

The established conclusions are already stronger than a collection of
isolated transcendence examples. The cofinite span determines the whole
generated finite-type algebra, the finite exceptional field, and the
geometric splitting of every relation. The same mechanism constructs a
continuum of positive infinitesimal surreal numbers with all derivatives
independent over a full rational-coefficient Hahn field, and a continuum
of common-domain surcomplex analytic functions with all mixed jets
independent. These are precise, nontrivial research results offered with
complete proofs and explicit limitations, while their bibliographic
novelty remains a matter for independent review.

\clearpage
\appendix
\section{Notation and scope at a glance}
\small
\begin{longtable}{@{}p{.23\textwidth}p{.71\textwidth}@{}}
\toprule
Symbol & Meaning and scope\\
\midrule
\endfirsthead
\toprule Symbol & Meaning and scope\\\midrule
\endhead
$M((t^\Gamma))$ & Hahn field: coefficients in $M$, well-ordered support in the set-sized ordered abelian group $\Gamma$.\\
$W$ & Cofinite span of coefficient vectors; intersection over finite deletions, not arbitrary ordinal tails.\\
$E$ & Finite set of vectors outside $W$; precisely the exceptional coefficient indices in the classification.\\
$B_E$ & $M(r_i:i\in E)((t^\Gamma))$; the exact algebraic part of $B(Y)$, of degree $2^{|E|}$ over $B$.\\
$\AR$ & Real-algebraic numbers, a countable subfield of $\R$.\\
$B_0\subset H_0$ & $\Q((t^\Q))\subset\AR((t^\Q))$, concretely embedded in $\No$ by $t=\omega^{-1}$.\\
$\BM$ & Normalized Berarducci--Mantova derivation on actual surreals; its restriction to $H_0$ is $-t\EDer$.\\
$\EDer$ & Euler derivation; $\EDer(t^q)=qt^q$ and coefficient-field elements are fixed.\\
$\partial$ & In analytic sections, ordinary differentiation in the independent parameter $z$ on coefficient germs.\\
$B$ in analytic sections & $\C(z)((t^\Q))$, a formal relation field, not a common-domain holomorphic ring.\\
$K_\Q$ & $\C((t^\Q))\subset\No[i]$; $K_\Q(z)$ is a smaller relation field contained in $B$.\\
$D$ & The ordinary complex disk $|z|<1$, shared by every analytic coefficient in every constructed function.\\
$\cf$ & The cardinal $2^{\aleph_0}$, not the proper-class size of all surreals.\\
\bottomrule
\end{longtable}
\normalsize

\phantomsection
\addcontentsline{toc}{section}{References}
\begin{thebibliography}{99}
\small\raggedright
\setlength{\itemsep}{3pt}
\bibitem{Repo}
Vladimir Reshetnikov, \emph{Surreal}, project README, revision
\texttt{\repocommit}, inspected September 22, 2026.
\href{https://github.com/VladimirReshetnikov/Surreal/blob/608dd23c0539fce73723843949ee627641c27b30/README.md}{Pinned project README}.

\bibitem{RepoMap}
Vladimir Reshetnikov, \emph{The surreal and surcomplex reports},
documentation catalogue, same revision and access date.
\href{https://github.com/VladimirReshetnikov/Surreal/blob/608dd23c0539fce73723843949ee627641c27b30/docs/README.md}{Pinned documentation catalogue}.

\bibitem{RepoDiffEq}
Vladimir Reshetnikov, \emph{Differential Algebra and Differential Equations
over the Surreal and Surcomplex Numbers}, companion report in this
collection, \texttt{docs/surcomplex/differential-equations/}.
Its section ``The rational-exponent workspace'' fixes
$\C((t^\Q))$ with $t=\omega^{-1}$ at \texttt{diff:eq:HahnQ}, and its
Proposition \texttt{diff:prop:hahnderiv} proves the derivative formula
\texttt{diff:eq:hahnderiv} on a real-exponent workspace. That report is
the statement of record for those two facts; they are cited, not reproved,
in Section~\ref{tail:sec:surreal}. Its warning
\texttt{diff:warn:monomialmap} separates the Conway monomial map from
general surreal exponentiation.

\bibitem{RepoDyn}
Vladimir Reshetnikov, \emph{Surcomplex Dynamics: Linearization Thresholds,
Periods, and Normal Forms}, companion report in this collection,
\texttt{docs/surcomplex/dynamics-and-normal-forms/}.
Its warning \texttt{dyn:warn:diffalg} states that the transcendence in its
algebraicity dichotomy \texttt{dyn:thm:dichotomy} must \emph{not} be
upgraded to differential transcendence, because its own solutions satisfy a
first-order linear equation over its base. That warning is true of those
solutions and is not contradicted here; see
Remark~\ref{tail:rem:dynupgrade}.

\bibitem{StacksGalois}
The Stacks Project Authors, \emph{Fields, Theorem 9.21.7: Fundamental theorem of Galois theory},
Tag 09DW, consulted September 22, 2026.
\url{https://stacks.math.columbia.edu/tag/09DW}.

\bibitem{StacksNormal}
The Stacks Project Authors, \emph{Fields, Lemma 9.15.7(2): extension of
embeddings in a normal algebraic extension}, Tag 0BME.
\url{https://stacks.math.columbia.edu/tag/0BME}.

\bibitem{StacksKummer}
The Stacks Project Authors, \emph{Fields, Section 9.24: Kummer extensions},
Tag 09I6, consulted September 22, 2026.
\url{https://stacks.math.columbia.edu/tag/09I6}.
Only the basic Kummer background is used; the multiquadratic statement
needed here is proved in Lemma~\ref{tail:lem:signs}.

\bibitem{vdH}
Joris van der Hoeven, \emph{Operators on generalized power series},
Illinois Journal of Mathematics \textbf{45} (2001), no.~4.
\url{https://doi.org/10.1215/ijm/1258138061}.

\bibitem{KKS}
Elliot Kaplan, Lothar Sebastian Krapp, and Michele Serra,
\emph{Decomposing the Automorphism group of the surreal numbers},
arXiv:2509.22374, consulted September 22, 2026; especially the
normal-form and Hahn-field preliminaries.
\url{https://arxiv.org/abs/2509.22374}.

\bibitem{BM}
Alessandro Berarducci and Vincenzo Mantova,
\emph{Surreal numbers, derivations and transseries},
Journal of the European Mathematical Society \textbf{20} (2018),
339--390. DOI: 10.4171/JEMS/769. Preprint arXiv:1503.00315.
\url{https://arxiv.org/abs/1503.00315}.

\bibitem{MM}
\v Zarko Mijajlovi\'c and Branko Male\v sevi\'c,
\emph{Differentially Transcendental Functions},
Bulletin of the Belgian Mathematical Society--Simon Stevin
\textbf{15} (2008), no.~2. Preprint arXiv:math/0412354.
\url{https://arxiv.org/abs/math/0412354}.
\end{thebibliography}
\end{document}
